%% !TEX root = manuscript.tex

\section*{Remerciements}
Je souhaite exprimer ma sincère gratitude à toutes les personnes qui ont contribué à mon
intégration et à ma progression professionnelle tout au long de mon alternance au sein de la
Holding SPA du groupe Angelotti.

En premier lieu, je remercie chaleureusement ma tutrice en entreprise, Sabine Cruanas, pour
son accompagnement constant, sa bienveillance et le partage de ses connaissances. Son
expertise et son soutien ont été déterminants dans ma montée en compétences, tant sur le
plan technique que méthodologique.

Je tiens également à remercier Mme Trottier, ma tutrice universitaire, pour son suivi
rigoureux et ses conseils avisés. Je remercie aussi l'ensemble des équipes d'Angelotti avec
lesquelles j'ai collaboré, en particulier la direction financière et les équipes
commerciales, dont la disponibilité et les retours constructifs ont été essentiels pour
adapter mes réalisations aux besoins concrets de l'entreprise. C'est en écoutant leurs
objections que la plupart de mes livrables ont pris leur forme définitive.

Enfin, je remercie le groupe Angelotti pour m'avoir offert cette opportunité d'alternance
dans un environnement professionnel enrichissant, ainsi que mon université et l'équipe
pédagogique du Master MIASHS pour leur accompagnement tout au long de ce parcours. Merci à
toutes et à tous pour votre confiance et votre soutien.

\clearpage
\section*{Résumé}
Ce mémoire présente les travaux que j'ai conduits durant ma deuxième année de Master MIASHS,
en alternance au sein de la Holding SPA du groupe Angelotti, promoteur-aménageur implanté en
Occitanie et en région PACA. L'année a été dominée par une transformation structurante du
système d'information : le remplacement de l'ERP historique Grimmo par la solution SPO,
éditée par Salvia.

Le premier semestre a été celui de l'ingénierie des données. J'y ai conçu et industrialisé
la reprise des données opérationnelles vers le nouvel ERP, sur un périmètre de 283
opérations immobilières, en m'appuyant sur les interfaces de programmation de l'éditeur. J'y
expose les difficultés algorithmiques rencontrées, notamment la gestion des identifiants
dynamiques et du versionnement des données budgétaires, ainsi que la stratégie de
déploiement en «~coquilles vides~» adoptée pour garantir la continuité de la facturation. En
parallèle, j'ai assuré le maintien du système décisionnel et participé au déploiement du
Single Sign-On sur deux périmètres du groupe.

Le second semestre a prolongé ce travail de structure, puis l'a valorisé. J'ai d'abord conduit
la reprise des tiers, c'est-à-dire des clients acquéreurs, en concevant une chaîne de traitement
entièrement pilotée par formules qui a permis de créer 1~434 tiers dans SPO à partir de 1~602
lignes brutes, sans doublon détecté par les trois clés d'unicité retenues. J'ai ensuite mené de
bout en bout un projet de Data Science, le «~Copilote Financier~», qui outille la direction
financière sur trois questions de gestion : le risque de dérive des marges, le rythme
d'écoulement des programmes et le pilotage des prix du stock.

Ce mémoire relit ensuite ces quatre missions transversalement, sous l'angle de la conduite de
projet, des technologies et des méthodes employées, des régimes de preuve mobilisés pour
évaluer chaque livrable et des difficultés rencontrées.

De cette année, il tire une thèse qu'il s'attache à démontrer plutôt qu'à énoncer : la
structure de la donnée commande ce que la modélisation peut en tirer. Il se termine par
une réflexion sur le métier de data scientist tel que je l'ai pratiqué, et sur le projet
professionnel qui en découle.
